\section{Solipsis Protocol Details}

\newcommand{\nd}{\ensuremath{\textsc{nd}}\xspace}
\newcommand{\UID}{\ensuremath{\textsc{UID}}\xspace}
\newcommand{\pos}{\ensuremath{\textsc{pos}}\xspace}
\newcommand{\sn}{\ensuremath{\textsc{sn}}\xspace}

\newcommand{\clist}{\ensuremath{\textsc{cList}}\xspace}
\newcommand{\descs}{\ensuremath{\textsc{descs}}\xspace}
\newcommand{\rn}{\ensuremath{\textsc{rNet}}\xspace}
\newcommand{\gn}{\ensuremath{\textsc{getNeighs}}\xspace}
\newcommand{\route}{\ensuremath{\textsc{route}}\xspace}
\newcommand{\cursn}{\ensuremath{\textsc{cursn}}\xspace}
\newcommand{\cupdate}{\ensuremath{\textsc{cUpd}}\xspace}
\newcommand{\aupdate}{\ensuremath{\textsc{aUpd}}\xspace}


\newcommand{\olist}{\ensuremath{\textsc{ownList}}\xspace}

\newcommand{\vn}{\ensuremath{\textsc{ver}}\xspace}

\newcommand{\msgentry}[1]{\vspace{2ex}\noindent{\large #1}\\}
\newcommand{\recby}[1]{\textbf{received by:} #1\\}
\newcommand{\param}[1]{\textbf{parameters:} #1\\}
\newcommand{\desc}[1]{\\}

\newcommand{\dstruct}[1]{\vspace{2ex}\noindent{\large #1}\\}
\newcommand{\mby}[1]{\textbf{maintained by:} #1\\}




%data structures

%renderer 
\newcommand{\modtab}{\ensuremath{\textsc{3Dmodels}}\xspace}
%table containing an entry for each object for which there is at least
%a cached 3D model. 

%avatar node
%\newcommand{\vlist}{\ensuremath{\textsc{vList}}\xspace}
\newcommand{\snode}{\ensuremath{\textsc{sNode}}\xspace}

%site node
\newcommand{\vlist}{\ensuremath{\textsc{vList}}\xspace}


%management of 3D descriptions
\newcommand{\modreq}{\ensuremath{\textsc{3Dreq}}\xspace}
\newcommand{\modrep}{\ensuremath{\textsc{3Dresp}}\xspace}


%management of views
\newcommand{\avhb}{\ensuremath{\textsc{avNodeHBeat}}\xspace}
%from avatar to site, periodic

\newcommand{\snhb}{\ensuremath{\textsc{siteNodeHBeat}}\xspace}
%from site to avatar, periodic

\newcommand{\vupd}{\ensuremath{\textsc{vUpdate}}\xspace}
%from site to site, periodic

\newcommand{\avupd}{\ensuremath{\textsc{avUpdate}}\xspace}
%from avatar to avatar, when avatar moves and with minimum frequency

\newcommand{\ncache}{\ensuremath{\textsc{newCopy}}\xspace}
%from avatar to avatar or site, when caching 3D description


%%%%%%%%%%%


\newcommand{\safed}{\ensuremath{\textsc{safeD}}\xspace}
\newcommand{\neighd}{\ensuremath{\textsc{neighD}}\xspace}





\subsection{Data Structures Maintained at Node}
Solipsis nodes maintain a set of data structures to manage the
interaction with the objects that constitute the 3D world. In the
following, we describe each of these structures for site nodes and for
avatar nodes, respectively.

%\subsubsection{Site Nodes}

\dstruct{\vlist}
\mby{Site Node}
Each site node maintains a visible list, \vlist. The \vlist is a table
containing an entry for each object that is visible from the current
site node's cell. Each entry in the list consists of an object
descriptor and of a timestamp indicating when the entry was last
updated.  Each object in the table may represent a site or an avatar
and may be inside or outside of the current node's cell. However, the
implementation may use separate tables for each of these object
categories, as long as the protocol's semantics remain the same as
specified in this document.

% \paragraph{Changed List}
% The changed list \clist is a complement to the visible list described
% above and records the identifiers of the objects in the visible list
% whose descriptors have been recently updated (for more details see
% Section~\ref{}).

% % A more fine grained distinction is
% % provided by the \emph{visible site list}, the \emph{avatar list}



\dstruct{\vlist}
\mby{Avatar Node}
Similar to site nodes, avatar nodes also maintain a visible list,
\vlist, containing an entry for each object that is potentially
visible from the avatar's location. As in the case of site nodes, each
object in the table may represent a site or an avatar in any of the
surrounding cells, although the implementation may also use a separate
table for each object category.


\dstruct{\snode}
\mby{Avatar Node}
Each avatar maintains an $\snode$ pointer, that always refer to site
node that is currently responsible for the avatar's location.

\dstruct{\modtab}
\mby{Host}
Each host maintains a model table \modtab that is accessible by all
avatar and site nodes associated with the host. The table contains
references to local copies of 3D models cached by the current host. At
each instant the \modtab table must always contain references to the
3D model of sites and avatars permanently associated with the
hosts. In addition, it may contain references to remote avatars and
sites as required for visualization purposes or in order to achieve
persistency in the presence of failures.

\newcommand{\achange}{\ensuremath{\textsc{avChange}}\xspace}
\newcommand{\adesc}{\ensuremath{\textsc{avDesc}}\xspace}
\newcommand{\newsnode}{\ensuremath{\textsc{NewSiteNode}}\xspace}
\newcommand{\cursite}{\ensuremath{\textsc{curSite}}\xspace}


\subsection{RayNet}
Both avatar and site nodes have access to a \rn object that abstracts
their interface with the RayNet overlay network. The \rn object
provides each site node with with the following two functions.
\begin{itemize}
\item $\gn()$: may be called only by site nodes and returns the list
  of the node's neighbors.
\item 
$\route(m, p)$: may be called by both avatars and site nodes and routes a
message $m$ to the site node associated with the location $p$ in the
coordinate space. 
\end{itemize}



\subsection{Protocol Messages}
Protocol Messages are exchanged asynchronously between Solipsis nodes
using UDP. This may be overridden when the two communicating nodes (e.g. a
site and an avatar) actually reside on the same peer. 

% \paragraph{\achange}
% \recby{site node}
% \param{set of (name, value) pairs} 
% Inform the receiving site node of a change in one of the 

\subsubsection{View synchronization between site and avatar nodes and
  between site nodes}
\msgentry{\avhb}
\recby{site node}
\param{ObjDescriptor nd}
\desc

Each avatar periodically sends an avatar heartbeat message to the site
node responsible for the avatar's current location, $\snode$. The
interval between messages, $T_{hb}$ should be on the order of
seconds. However, avatar nodes may occasionally send more frequent
updates, if their position or other features have undergone
significant changes. The receiving site node reacts to the an \avhb
message by updating its \vlist. Specifically, the entry associated
with the avatar is updated either if it was absent from the \vlist or
if the received entry has a fresher sequence number than the one in
the table.


\msgentry{\vupd}
\recby{site node, avatar node}
\param{list$<$ObjDescriptor$>$ nds}
\desc

A view update message contains the necessary information to update the
visible lists neighboring site nodes. The message is sent every $T_u$
seconds by every site node to each of its neighboring site nodes
($\rn.\gn()$). The contents of the message consist of a list of node
descriptors representing all the objects from the sender node's
visible list that are also visible from the sender node's current
position. A site node that receives a \vupd message from one of its
neighbors in the raynet overlay updates its visible list with all the
node descriptors in the message.

\msgentry{\snhb}
\recby{avatar node}
\param{long \sn, list$<$ObjDescriptor$>$ nds}
\desc

A view update message contains the necessary information to update the
visible lists of the avatar nodes in a site node's cell. Similar to
\vupd messages, \snhb messages are sentevery $T_u$ seconds by each
site node to all of the avatars that are currently in their cells. In
this first version of the protocol sending should be done directly to
all avatars in the cell. Future versions will include a more efficient
data distribution strategy to reduce the forwarding load on site
nodes.  

Different from the case of \vupd, the contents of the a \snhb
message comprise all the node descriptors of the objects that are in
the sender site node's visible list. An avatar receiving the message,
uses it to update its visible list.  Specifically it adds to it all
node descriptors in nds which are not already present. In addition, it
updates each entry in the list for which there is a newer node
descriptor (with a fresher sequence number) in nds.

\emph{It should be noted that $T_{u}$ may be on the order of seconds
  or tens of seconds.}

%The pseudocode fragments corresponding
%to these interpretations are given in Algorithm~\ref{alg:vupdate}

% \begin{algorithm}
% \begin{algorithmic}

% \Procedure{sendChangedUpdate}{}{Executed every $T_c$ seconds by each site
% node.}
% \State \cupdate.\sn=\cursn++
% \State \cupdate.\descs=\clist
% \State $\clist=\emptyset$
% \ForAll {$n \in \rn.\gn$}
%   \State send \cupdate to $n$
% \EndFor
% \ForAll {$d \in \olist$}
%   \State send \cupdate to $d.\UID$
% \EndFor
% \EndProcedure
% \end{algorithmic}
% \caption{site node sending \cupdate and \vupd messages}
% \label{alg:scvupdate}
% \end{algorithm}



% \begin{algorithm}
% \begin{algorithmic}
% \Procedure{sendViewUpdate}{}{Executed every $T_v$ seconds by each site
% node.}
% \State \vupd.\sn=\cursn++
% \ForAll {$\nd \in \vlist$}
%    \State \nd.\sn=\nd.\sn++
%    \If {isInMyCell(\nd.\pos)}
%       \olist[\nd.\UID]=\nd
%    \EndIf
% \EndFor
% \State \vupd.\descs=\vlist
% \ForAll {$n \in \rn.\gn$}
%   \State send \vupd to $n$
% \EndFor
% \ForAll {$d \in \olist$}
%   \State send \vupd to $d.\UID$
% \EndFor
% \EndProcedure

% \Procedure{receiveViewUpdate}{\vupd msg}{Upon Receipt of a \vupd
%   message by each avatar or site node.}
% \ForAll {$\nd \in msg.\descs$}
%   \If {\vlist contains \nd.\UID}
%      \If {\vlist.get(\nd.\UID).\sn $$<$$ \nd.\sn}
%         \State \vlist[\nd.\UID]=\nd
%      \EndIf
%   \Else
%      \State \vlist[\nd.\UID]=\nd
%   \EndIf
% \EndFor
% \EndProcedure



\subsubsection{View synchronization between avatar nodes}


\msgentry{\avupd}
\recby{avatar node}
\param{ObjDescriptor nd}
\param{set$<$ObjDescriptor$>$ neighborhood}
\desc
Avatar update messages are exchanged between neighboring avatars to
maintain up-to-date information about close objects and particularly
about those that may interact with their movements. Each avatar node
should send an \avupd message whenever it changes its position or
other attributes such as parameters of its current
animation. Moreover, it should send additional \avupd messages to
prevent their frequency from falling below one update every $T_{upd}$
seconds.

Each avatar node should compute the set of recipients of avatar update
messages based on its size and that of the avatar nodes in its cell as
specified in the node descriptors. Specifically, the set of recipients
should include all the avatar that are at a distance that is less than
or equal to the sum of their vRadius values, plus the sender avatar's radius,
plus a safe distance parameter $\safed$. This guarantees that each
update reaches all the avatars that may potentially collide with the
current avatar before the next update\footnote{During standard
  operation}.

Each \avupd message comprises two parts. The first is the
node descriptor of the avatar sending the message. The second is a
list of node descriptors of nodes from the sending avatar's \vlist,
which includes all the avatars that are currently closer to the
sending avatar than to the recipient of the message and that are at a
distance that is less than $\neighd$ from the latter. \df{To get
  better guarantees, we need a raynet-like structures for avatars as
  well.}

An avatar node reacts to the receipt of an \avupd by updating its
visible list. Specifically, it inserts all node descriptors in the
message that are not already present in the list. Moreover, it updates
any entry for which the sequence number associated with the node
descriptor in the \avupd message is fresher than the sequence number
associated with the node descriptor that is currently in the visible
list.


\subsubsection{Managing 3D Descriptions}


\msgentry{\ncache}
%\sentby{host}
\recby{site node, avatar node}
\param{UID, \vn, $f$, url}
\desc
The \ncache message informs an avatar or a site node that a host has
cached a new copy of \vn version of the file $f$ associated with the
object identified by UID. 
The receiving node updates its private copy of the object descriptor
to include a reference to the new cached copy. \df{Need to agree on
  what to do with old versions}



\msgentry{\modreq}
\recby{host}
\desc
Hosts may obtain copies of 3D models by contacting one or more nodes
that have cached their latest version as specified in object
descriptors. To achieve this, they should send a \modreq message
specifying  the identifier of the object for which they are retrieving
the model, the type of model file they are retrieving (e.g. wireframe,
animation, ...), and the desired version number. 
The message is addressed to hosts regardless of whether they host site
or avatar nodes, and it triggers a response consisting of a \modrep
message as specified in the following.


\msgentry{\modrep}
\desc
A host receiving a \modreq message from another host should react by
looking up the desired 3D model in its own model table. If the table
contains a version of the required 3D model that is at least as recent
as the one requested by the \modreq, then it responds with a \modrep
message that includes a copy of such version. 



% \Procedure{Avatar.receiveViewUpdate}{\vupd msg}{Upon Receipt of a \vupd
%   message by an avatar node.}
% $\vlist=\emptyset$
% \If {msg.\sn $>$ sitesSN}
%   sitesSN=msg.\sn
%   \ForAll {\nd \in msg}
%     \State \vlist[\nd.\UID]=\nd
%   \EndFor
%   \EndIf
% \EndProcedure

% \end{algorithmic}
% \caption{site node sending and site or avatar nodes receiving
%   \vupd messages}
% \label{alg:vupdate}
% \end{algorithm}


% \msgentry{\cupdate}
% \recby{site node, avatar node}
% \param{list$<$ObjDescriptor$>$ nds}
% A change update message is structurally identical to a view update
% message and performs a similar task: updating information about
% visible objects. However, different from \vupd messages, \cupdate
% messages only contain information about the objects whose descriptors
% have changed since the last \cupdate or \vupd and are sent by site
% nodes whenever they detect a change in the set of objects in their
% visible list. To avoid flooding the network with \cupdate messages in
% the case of a very dynamic 3D world, nodes SHOULD avoid sending more
% \cupdate messages than one every $T_c$ seconds. If a node detects a
% change less than $T_c$ seconds after the last \cupdate, it SHOULD
% defer the new \cupdate so as not to exceed the maximum frequency of
% one every $T_c$ seconds. A reasonable value for $T_c$ should be on the
% order of $0.1 T_v$. The pseudocode fragment triggered by the receipt
% of a \cupdate message from a site or an avatar node are given in
% Algorithm~\ref{algo:reccupdate}.


% \begin{algorithm}
% \begin{algorithmic}
% \Procedure{sendChangeUpdate}{}{Executed by each site node every $T_c$
%   seconds}
% \If {$\clist \neq \emptyset$}
%    \State \cupdate.\descs=\clist
%    \State $\clist=\emptyset$
%    \ForAll {$n \in \rn.\gn$}
%      \State send \cupdate to $n$
%    \EndFor
%    \ForAll {$d \in \olist$}
%      \State send \cupdate to $d.\UID$
%    \EndFor
% \EndIf
% \EndProcedure


% % \EndIf
% % \ForAll {$\nd \in \vlist$}
% %    \State \nd.\sn=\nd.\sn++
% %    \If {isInMyCell(\nd.\location)}
% %       \olist[\nd.\UID]=\nd
% %    \EndIf
% % \EndFor
% % \State \vupd.\descs=\vlist
% % \ForAll {$n \in \rn.\gn$}
% %   \State send \vupd to $n$
% % \EndFor
% % \ForAll {$d \in \olist$}
% %   \State send \vupd to $d.\UID$
% % \EndFor
% % \EndProcedure

% \Procedure{receiveViewUpdate}{\vupd msg}{Upon Receipt of a \vupd
%   message by each avatar or site node.}
% \ForAll {$\nd \in msg.\descs$}
%   \If {\vlist contains \nd.\UID}
%      \If {\vlist.get(\nd.\UID).\sn $$<$$ \nd.\sn}
%         \State \vlist[\nd.\UID]=\nd
%      \EndIf
%   \Else
%      \State \vlist[\nd.\UID]=\nd
%   \EndIf
% \EndFor
% \EndProcedure


% % \Procedure{Avatar.receiveViewUpdate}{\vupd msg}{Upon Receipt of a \vupd
% %   message by an avatar node.}
% % $\vlist=\emptyset$
% % \If {msg.\sn $>$ sitesSN}
% %   sitesSN=msg.\sn
% %   \ForAll {\nd \in msg}
% %     \State \vlist[\nd.\UID]=\nd
% %   \EndFor
% %   \EndIf
% % \EndProcedure

% \end{algorithmic}
% \caption{site node sending and site or avatar nodes receiving
%   \vupd messages}
% \label{alg:vupdate}
% \end{algorithm}


% \msgentry{\aupdate}
% \recby{site node}
% \param{long sequenceNumber}
% \param{ObjectDescriptor sd} 
% \aupdate messages are the means used by avatar nodes to modify the
% object descriptor for their avatar: either to update the avatar's
% physical characteristics or ( identified by \nd.\UID) to move it to a
% different location. 

% Each avatar SHOULD route an \adesc message to the site node
% responsible for its current location every $T_u$ seconds. In addition,
% it MAY also inform the site node whenever the user attempts to move
% the avatar. As in the case of \cupdate messages, however, the
% frequency of \adesc messages from a given avatar node MUST not exceed
% one every $T_c$ seconds. Avatar nodes MUST therefore defer or cancel
% \aupdate messages that would cause this frequency limit to be
% exceeded.

% Site nodes respond to the receipt of an \aupdate message by updating
% the node descriptor in their visible lists and updating the
% corresponding sequence number, as shown in
% Algorithm~\ref{alg:recaupdate}.
















% \section{Protocol Pseudocode}
% This section of the specification describes the operations carried out
% by site and avatar nodes in terms of pseudocode. The pieces of
% pseudocode should be interpreted as guidelines for the implementation
% of the protocol, and may be modified as long as its semantics are not
% changed.

% \subsubsection{Periodic Actions}
% Both site and avatar nodes carry out periodic actions to maintain
% consistent information in their visibility lists. The following is a
% summary of these actions.

% \msgentry{Site Nodes}
% \begin{itemize}
% \item send a \vupd message to every neighbor every $T_v$ seconds.
% \item send a \vupd message to every avatar in own cell every $T_v$ seconds.
% \end{itemize}
% \msgentry{Avatar Nodes}
% \begin{itemize}
% \item send a \apropose message to the site responsible for their
%   location (using the RayNet's route primitive) every $T_v$ seconds.
% \end{itemize}
% \subsubsection{Exogenous Actions}
% Exogenous actions comprise all the interactions of the protocol with
% the outside world, that is with the Solipsis User Interface, as well
% as reactions to unexpected disconnections or network failures.

% \msgentry{Avatar Moves}
% The most basic form of interaction between a user and the \sol world
% is the movement of controlled avatars. Each user should be able to
% control the movements and the other aspects of the behavior of its
% avatar, ultimately determining its interaction with the \sol
% world. The movement of avatars however, is also subject to the laws of
% physics and so is their interaction with other objects in the \sol
% world, including other avatars. To achieve consistency between the
% commands of users and the constraints dictated by the environment, we
% have each site node act as a coordinator for the movements of the
% avatars that are currently in its cell.\df{Will this scale properly?
%   Do we need a more complex protocol in which avatar nodes also take
%   an active part in collision detection and physics computations?}

% Avatar nodes interact with site nodes by proposing movements. Site
% nodes react by computing the resulting scene and communicating it to
% the avatars. \df{Ultimately, this is what turned out to make the most
%   sense to me. Please let me know if you agree with this!}



% An avatar node reacts to a change in the location of its avatar by
% sending an \adesc message to its current site node, that is the one
% associated with  location before the move.

% \begin{algorithm}
% \caption{Avatar moves from one location $l_0$ to another $l_1$}
% \begin{algorithmic}
% \Procedure{avatarMove}{$l_0$, $l_1$}{}
% % \State \desc.\pos=\l
% % \State $n_0$=\rn.\gn($l_0$)
% % \State $n_1$=\rn.\gn($l_1$)
% % \If {$n_0 \neq n_1$}
% \desc.\pos=$l_1$
% \rn.\route(\desc, $l_0$)
% \EndProcedure
% \end{algorithmic}
% \end{algorithm}





% A site node receiving an \adesc message performs the following steps.

% \begin{algorithm}
% \caption{site node receiving an \adesc message}
% \begin{algorithmic}
% \Procedure{receiveAvatarDescriptorFrom}{\adesc}{}
% \If {isVisibleFromMyCell(\nd.\vis, \nd.\pos)}
%   \If {\vlist contains \nd.\UID}
%      \If {\vlist.get(\nd.\UID).\sn $$<$$ \nd.\sn}
%         \State \vlist[\nd.\UID]=\nd
%      \EndIf
%   \Else
%      \State \vlist[\nd.\UID]=\nd
%   \EndIf
% \Else
%   \If {\vlist contains \nd.\UID}
%      \State \vlist.remove(\nd.\UID,\nd)
%   \EndIf
% \EndIf

% \If {isInMyCell(\nd.\pos)
%   \State \olist[\nd.\UID]=\nd
% \Else
%   \State \olist.remove(\nd.\UID)
% \EndIf
% \EndProcedure

% \Function {isVisibleFromMyCell}{$r$, $c$}{returns true if the current
%   site node's cell intersects the visibility region defined by the
%   circle with radius  $r$ centered at $p$}
% TBD\\
% \EndFunction

% \Function {existsNeighborInRegion}{$r$, $c$}{returns true if there is
%   a neighbor in the circular region centered at $p$ with radius $r$}
% \ForAll {n \in \rn.\gn}
%   \If (n.\pos \in \region($c$, $r$))
%     \Return \True
% \EndFor
% \Return \False
% \EndFunction

% \Function {existsNeighborInRegion}{$r$, $c$}{returns true if there is
%   a neighbor in the circular region centered at $p$ with radius $r$}
% \ForAll {n \in \rn.\gn}
%   \If (n.\pos \in \region($c$, $r$))
%     \Return \True
% \EndFor
% \Return \False
% \EndFunction

% \Function {isInMyCell}{$p$}{returns true if position $p$ is in the
%   current site node's cell}
% myD=\distance(n.\pos, \desc.\pos)
% \ForAll {n \in \rn.\gn}
%   \If {\distance(n.\pos, $p$)$>$myD}
%     \Return \False
% \EndFor
% \Return \True
% \EndFunction
% \end{algorithmic}
% \end{algorithm}

% \msgentry{View Update and Change Update Messages}
% Periodically, site nodes inform the avatars in their cells as well as
% neighboring site nodes about visible objects. This is achieved through
% the combination of two mechanisms. First, site nodes send a \vupd
% message every $T_v$ seconds to neighboring sites and avatars. This
% message contains copies of the descriptors of all the objects visible
% from the site node's cell. Second, with a frequency of up to one every
% $T_c$, where $T_c$<$$<$ T_v$, site nodes also send \cupdate messages to
% inform neighboring sites and visiting avatars about changes occurred
% from the previous update message. As described above, \cupdate
% messages also consists of a list of node descriptors. However, while
% the absence of a descriptor from a \vupd message implies that the
% corresponding object is no longer within visibility of the sender
% node's cell, its absence from a \cupdate message simply implies that
% the object's descriptor has not changed since the last update.
% \inlinedf{Right now all updates are sent via unicast to site and
%   avatar nodes. This is probably fine for messages addressed to site
%   nodes. Messages addressed to avatar nodes, however, may benefit from
%   the use of some kind of multicast tree.}

% \begin{algorithm}
% \caption{site node receiving an \adesc message}
% \begin{algorithmic}
% \Procedure{sendViewUpdate}{}{Executed every $T_v$ seconds by each site
% node.}
% \State \vupd.\sn=\cursn++
% \State \vupd.\descs=\vlist
% \ForAll {$n$ \in \rn.\gn}
%   \State send \vupd to $n$
% \EndFor
% \ForAll {$d$ \in \olist}
%   \State send \vupd to $d.\UID$
% \EndFor
% \EndProcedure

% \Procedure{sendChangedUpdate}{}{Executed every $T_c$ seconds by each site
% node.}
% \State \cupdate.\sn=\cursn++
% \State \cupdate.\descs=\clist
% \State $\clist=\emptyset$
% \ForAll {$n$ \in \rn.\gn}
%   \State send \cupdate to $n$
% \EndFor
% \ForAll {$d$ \in \olist}
%   \State send \cupdate to $d.\UID$
% \EndFor
% \EndProcedure

% \end{algorithmic}
% \end{algorithm}



% \item It verifies whether the avatar's location as specified in the
%   node descriptor is inside its own cell. 
%   If this is NOT the case then it uses the RayNet to forward the 


% If the site node's visible map does not contain an entry for \nd.\UID,
% then a new entry should be created and associated with the received
% descriptor \nd. If, instead, the map already contains such an entry,
% then the it SHOULD be update only if the if the version number
% nd.version is greater than the version number associated with the
% entry in the table. Otherwise the message MUST be silently discarded.



% \msgentry{\newsnode}
% \recby{avatar node}
% \param{\cursite}
% Informs the avatar node that the site node responsible for its current
% location has changed. The avatar node MUST therefore\df{this and the
%   following could be SHOULD.} update its own \cursite reference to
% equal the one provided by the message.



%\msgentry{\avmoved}

% \msgentry{View Update and Change Update Messages}
% Periodically, site nodes inform the avatars in their cells as well as
% neighboring site nodes about visible objects. This is achieved through
% the combination of two mechanisms. First, site nodes send a \vupd
% message every $T_v$ seconds to neighboring sites and avatars. This
% message contains copies of the descriptors of all the objects visible
% from the site node's cell. Second, with a frequency of up to one every
% $T_c$, where $T_c$<$$<$ T_v$, site nodes also send \cupdate messages to
% inform neighboring sites and visiting avatars about changes occurred
% from the previous update message. As described above, \cupdate
% messages also consists of a list of node descriptors. However, while
% the absence of a descriptor from a \vupd message implies that the
% corresponding object is no longer within visibility of the sender
% node's cell, its absence from a \cupdate message simply implies that
% the object's descriptor has not changed since the last update.
% \inlinedf{Right now all updates are sent via unicast to site and
%   avatar nodes. This is probably fine for messages addressed to site
%   nodes. Messages addressed to avatar nodes, however, may benefit from
%   the use of some kind of multicast tree.}



%%% Local Variables: 
%%% mode: latex
%%% TeX-master: "specification"
%%% End: 

